\section{Introduction}\label{sec:intro}
Solvent extraction leaves oilseed meal wet with hexane. In a
desolventizer-toaster (DT), live steam strips the solvent while heated trays
supply additional energy \citep{sipos1961dt,witte_nd_soybeanmeal,karnofsky1985recovering}.
Residual solvent matters for safety and emissions; condensed water must
later be removed, and the meal's thermal history affects protein
quality \citep{epa1995ap42,kemper2013solvent,mustakas1981critical}.
The engineering question is how to supply enough heat and residence time
for solvent removal and toasting without excessive steam use, later
drying duty or damage to meal quality \citep{kemper2013solvent,witte_nd_soybeanmeal}.
A particle model must therefore connect solvent removal, moisture uptake
and heat transfer throughout the transition from externally wetted meal
to a particle containing only retained material and pore gas.

These processes are coupled. Evaporation consumes energy at the particle
surface or an internal liquid boundary. Incoming steam can condense, and
both retained water and residual hexane contribute to storage and transport.
Liquid-water imbibition also differs from vapour sorption. Experiments on soy
protein isolates show that imbibition combines sorption, capillarity and
interstitial retention, whereas vapour sorption describes a different
hydration experiment \citep{jovanovich2003wateruptake}.
Those measurements motivate treating liquid contact separately from vapour
exposure; their capacities and rates cannot be transferred directly
to hot, hexane-bearing extracted meal.

Published desolventization models emphasize different parts of this
sequence. The superheated-hexane experiments and receding-front model of
\citet{faner2019kinetics} describe a heat-fed liquid-hexane core. The transport
account of \citet{cardarelli2002modeling} resolves a diffusion-controlled
drying tail. \citet{coletto2022modeling} couples a transient radial particle
to a steady axial bed in the diffusion-controlled zone, treating the
preceding flashing zone separately. Here we combine component and energy
conservation at a receding liquid-hexane interface with binary pore gas,
matrix retention and a separately conserved external water film.
After the core disappears, the same stored components and energy continue
to evolve in a stationary radial particle.

In the first of two prescribed baths, a water-bearing particle in pure
superheated hexane is compared with the soybean experiment of
\citet{faner2019kinetics}; pure supplied vapour does not imply a
water-free meal. In the second, a particle is followed from feed
through a prescribed steam-rich DT history. Prescribing the bath isolates
the particle response, so the surrounding vapour composition and
temperature need not be predicted. In a future coupled tray calculation,
the signed water, hexane and energy fluxes would enter the gas and tray
balances, which would determine the next local bath state.

The two comparisons rest on different kinds of evidence. In the Faner
experiment, the measured constant-rate duty informs external heat
supply; later loading and temperature readings provide the comparisons.
Faner reports preheating in the apparatus and particle-size summaries, but
not the exact initial temperature, initial meal moisture or a mass-weighted
size histogram. The DT history is an engineering scenario; published outlet
ranges provide process context rather than measurements of the same
particle and operating point. Liquid-water imbibition is part of the particle
formulation, and its transport coefficients are declared rather than
measured. Numerical verification and complete time integration make
these assumptions inspectable; independent measurements are needed to
establish their predictive accuracy.

The paper has two goals: to state one conservative particle formulation
and to exercise it in both baths.
Section~\ref{sec:model} defines the particle equations and phase regimes,
and Section~\ref{sec:numerics} their numerical integration, admissibility
checks and verification. Section~\ref{sec:validation} presents the two
experiments, Section~\ref{sec:discussion}
discusses their meaning for DT practice and the measurements needed next,
and Section~\ref{sec:conclusions} concludes. The
Supplementary Material contains detailed source readings, parameters,
numerical evidence and the separate scalar Fickian comparison with the
earlier Faner measurements.
